% v005 stage C draft (2026-11-17): positioning of P3Net in the P3/GOMEA line. Replaces the v003 paragraphs
% "Linkage learning and P3", "Surrogate-assisted linkage learning" and "Closest prior work"
% (chapters/archive/v003/related_work/main.tex); the v003 pyramid figure (fig:p3-pyramid) can be reused unchanged.
% Sources: notes/ga_landscape_2026.md §16-§16d (literature search 2026-11-16/17). Dahi et al. 2026 is described
% from its supplementary material only (main text not accessible). The sentence on the surrogate result follows the
% current v005 + E1 results and is to be revised if the surrogate study continues.

\paragraph{Linkage learning and P3.} GOMEA and P3 build a linkage tree from statistical dependencies in the
population and use it to guide crossover, so that groups of dependent variables are not
disrupted~\citep{thierens2011optimal,goldman2014parameterless}. P3 replaces a single fixed-size population with a
pyramid of populations (Figure~\ref{fig:p3-pyramid}), adding a level only once existing levels stop improving, so that no population size is
chosen in advance; its evaluation complexity grows more slowly with problem size than that of five
competitors~\citep{goldman2015fast}. Multi-objective variants exist for both: MO-GOMEA~\citep{luong2014mogomea}
and MO-P3, applied to industrial process planning~\citep{przewozniczek2021mop3} and later extended with a Pareto-front
improvement phase~\citep{niemczyk2025pareto}. Unlike the gray-box setting that this line often
exploits~\citep{whitley2016graybox,bouter2023library}, NAS and HPO objectives are not decomposable into cheaply
recomputable subfunctions, so every variation step must be verified by a full evaluation or screened by a surrogate.

\paragraph{Linkage learning in NAS.} Members of this family have been applied to NAS before. den Ottelander et al.\
find MO-GOMEA, NSGA-II, local search and random search indistinguishable on NAS-Bench-101, with MO-GOMEA ahead of
random search on their larger MacroNAS spaces~\citep{denottelander2021localsearch}. ENCAS uses MO-GOMEA to search
cascades of networks drawn from supernetworks~\citep{chebykin2022encas}; Tran et al.\ apply GOMEA to NAS with a
training-free metric in place of training~\citep{tran2023gomeanas}; and Dahi et al., after applying linkage-learning
EAs to architecture search for quantum-circuit transpilation~\citep{dahi2024evolutionary}, use P3 to choose the depth
and layer widths of a small multilayer perceptron within a quantum-optimisation pipeline, a single-objective
architecture search~\citep{dahi2026tightening}. Bartnik adapts the surrogate-assisted P3-GOMEA
of~\citet{dushatskiy2021novelsurrogateassistedevolutionaryalgorithm} to multi-objective NAS on NAS-Bench-201
(accuracy against GPU energy), with an absolute regressor over the genotype as surrogate, an architecture-only
genotype, and MO-LS and surrogate-free MO-P3-GOMEA as baselines~\citep{bartnik2026evolutionary}; we include its two
variants, MO-DSA and MO-PSA, among our comparison arms.

\paragraph{Surrogates and linkage.} CS-GOMEA pairs GOMEA with a convolutional surrogate and names as open the
surrogate's failure to use linkage information in its own construction~\citep{dushatskiy2019csgomea}. LyMPuS builds
the surrogate from the discovered linkage itself and validates it with P3 on NK landscapes, Ising spin glasses and
Max3Sat; its empirical variant, eLyMPuS, discovers the dependency structure incrementally rather than assuming it
known~\citep{przewozniczek2026lympus}, and related work constructs monotonic surrogates from recursive linkage
discovery~\citep{przewozniczek2026monotonical}. P3Net's surrogate follows eLyMPuS in being relative and
linkage-aware over an incompletely discovered structure, but regresses a continuous fitness difference between
two solutions instead of making a discrete better/worse/ambiguous comparison, which places it closer to CS-GOMEA's
learned model and gives up eLyMPuS's guarantee of recovering the comparison under monotonicity.

\paragraph{Position of this study.} Against this background, P3Net is not the first use of linkage learning, or of
P3, in architecture search. What we have not found is P3 applied to the \emph{joint} multi-objective search over
architecture and training hyperparameters, nor a comparison of a linkage-learning optimiser with the main
Bayesian-optimisation, multi-fidelity and evolutionary methods of multi-objective NAS and HPO under one
pre-registered protocol that also reports the optimiser's own cost. This study supplies that comparison, with P3
in three forms (canonical, with a pairwise surrogate, and with eLyMPuS), on four primary spaces and five further
ones, and it reports where the surrogate does and does not change the outcome. All benchmarks are tabular or
surrogate-based (Section~\ref{sec:setup}).
